\chapter{Esperança que não pode ser hackeada}
\noindent\textit{Vinheta.} Na terceira noite do blackout, sem Wi-Fi e sem luz, uma igreja pequena se reúne no pátio. Um violão, vozes cansadas, crianças correndo. Alguém lê Ap 21: ``\emph{Eis que faço novas todas as coisas}''. Não mudou a rede; mudou o povo. Eles voltam para casa mais leves.

\section{Mapa bíblico (o eixo da esperança)}
\begin{itemize}[leftmargin=*,noitemsep]
  \item \textbf{O Cordeiro venceu} (Ap~5): a história não está sem autor.
  \item \textbf{Juízos com propósito} (Ap~8--16): revelar, corrigir, salvar (Rm~2:4--5).
  \item \textbf{Queda da Babilônia} (Ap~18): ídolos caem; Cristo permanece.
  \item \textbf{Nova criação} (Ap~21--22): Deus com o seu povo; fim da lágrima e da morte.
\end{itemize}

\section{Esperança cristã \emph{vs.} otimismo tecnológico}
Esperança não é \textbf{otimismo}. Otimismo diz: ``vai dar certo porque a curva de inovação sobe''. Esperança diz: ``\textit{Cristo reina}, ainda que a curva desça''. O primeiro depende de \emph{tendências}; o segundo, de \emph{promessas}.

\section{Como a esperança opera na prática}
\begin{itemize}[leftmargin=*,noitemsep]
  \item \textbf{Afeto}: aprende a \emph{desejar} a cidade por vir (Hb~11:10).
  \item \textbf{Ação}: trabalha aqui e agora (Jr~29:7; 1~Co~15:58).
  \item \textbf{Âncora}: volta-se sempre ao Cordeiro (Heb~6:19; Ap~5).
\end{itemize}

\section{Mini-exegese: juízo não é pânico}
Os selos, trombetas e taças são dramáticos, mas não aleatórios. Eles desmascaram ídolos e chamam ao arrependimento. Ler Apocalipse com esperança significa \textbf{reconhecer a justiça de Deus} e \textbf{abraçar sua paciência} (2~Pe~3:9), sem cinismo nem desespero.

\begin{nicbox}
\textbf{NIC da esperança}\\
\textbf{Narrativa}: a história é do Cordeiro.\\
\textbf{Identidade}: somos selados (Ap~7) --- não perfis descartáveis.\\
\textbf{Controle}: mesmo quando portas se fecham, Cristo abre a porta da perseverança (Ap~3:8).
\end{nicbox}

\section{Práticas que alimentam esperança}
\begin{checklist}
\begin{itemize}[leftmargin=*,noitemsep]
  \item \textbf{Salmos na mesa} (5 min): ler um salmo por refeição principal em voz alta.
  \item \textbf{Boletim de gratidão} (1/semana): a família lista 3 sinais do cuidado de Deus.
  \item \textbf{Visita curta} (quinzenal): encorajar alguém aflito com Ap~21:1--5.
  \item \textbf{Ofertar tempo} (mensal): servir alguém sem reciprocidade possível.
\end{itemize}
\end{checklist}

\section{Perguntas para grupos pequenos}
\begin{itemize}[leftmargin=*,noitemsep]
  \item Que \textit{otimismo} você precisa confessar como \emph{idolatria} tecnológica?
  \item Em qual promessa bíblica você precisa se ancorar nesta estação?
  \item Como sua igreja pode encarnar Ap~21 em pequenos gestos na cidade?
\end{itemize}

\bigskip
\noindent\textit{Próximo capítulo: A imagem que fala --- IA, culto e a idolatria da conveniência.}
